\documentclass[11pt,letterpaper]{article}
\usepackage{amsmath,amssymb,amsthm}
\newtheorem{proposition}{Proposition}
\newtheorem{corollary}[proposition]{Corollary}
\title{Elementary Facts About Divisibility}
\author{Sunniva Aldous}
\date{November 5, 2020}
\begin{document}
\maketitle

\section{Preliminaries}
Throughout, all letters denote integers unless stated otherwise. We write
$a \mid b$ if there is an integer $c$ with $b = ac$. The set of positive
integers is denoted $\mathbb{N}$, and $\mathbb{Z}$ is the set of all integers.

\begin{proposition}\label{prop:sum}
If $a \mid b$ and $a \mid c$, then $a \mid (b + c)$.
\end{proposition}
\begin{proof}
Write $b = ax$ and $c = ay$. Then $b + c = a(x + y)$.
\end{proof}

\begin{corollary}\label{cor:lin}
If $a \mid b$ and $a \mid c$, then $a \mid (mb + nc)$ for all $m, n \in
\mathbb{Z}$.
\end{corollary}

\section{Counting Divisors}
Let $d(n)$ denote the number of positive divisors of $n$. If
$n = p_1^{e_1} \cdots p_k^{e_k}$ is the prime factorisation, then
\begin{align}
d(n) &= \prod_{i=1}^{k} (e_i + 1), \label{eq:d}\\
\sigma(n) &= \prod_{i=1}^{k} \frac{p_i^{e_i+1} - 1}{p_i - 1}. \label{eq:sigma}
\end{align}
For instance, $d(360) = 4 \cdot 3 \cdot 2 = 24$, using \eqref{eq:d} and the
factorisation $360 = 2^3 \cdot 3^2 \cdot 5$.

\section{A Summation Identity}
Proposition~\ref{prop:sum} is used in the proof of the following identity,
which appears on page~\pageref{eq:id}:
\begin{equation}\label{eq:id}
\sum_{k=1}^{n} \left\lfloor \frac{n}{k} \right\rfloor = \sum_{m=1}^{n} d(m).
\end{equation}
Both sides count the pairs $(k, j)$ with $1 \le k j \le n$.

\subsection{Table of Small Values}
\begin{center}
\begin{tabular}{c|cccccc}
$n$ & 1 & 2 & 3 & 4 & 5 & 6 \\ \hline
$d(n)$ & 1 & 2 & 2 & 3 & 2 & 4 \\
\end{tabular}
\end{center}
\end{document}
